\section{Conclusion}
\label{sec:conclusion}

Respecting the published strategy constraints and solving the downstream games globally changes several conclusions in \cite{ShyStenbacka2005}.  In the Cournot model, the competition comparative static in equation (14) and Proposition 3 has the opposite sign from the one reported in the paper.  The all-active continuation is also incomplete after asymmetric sourcing histories.  Once zero-output regimes are included, Proposition 5 is no longer a global strategic-substitutes result, and the source duopoly can have an asymmetric pair of pure Stage-I equilibria alongside the symmetric equilibrium or, at a knife edge, a continuum.

The Hotelling model has a distinct issue.  The literal price game admits multiple pure corner continuations after sufficiently asymmetric sourcing histories, so the published interior backward induction is not globally single-valued.  Under a separate, explicit no-loss restriction \(p_j\ge c_j\), the pure price continuation is unique and the published symmetric sourcing candidate fails on the exact region derived in Proposition~\ref{prop:no-loss}.  That result is conditional on the auxiliary strategy restriction and is not a refinement claim about the literal source game.

The broader lesson is methodological rather than universal.  Two distinct corrections are involved.  The Cournot competition comparative static changes sign because the published symmetric expression is differentiated incorrectly; the sourcing cap then makes the corrected comparison weak or branch dependent when it binds.  Separately, nonnegative downstream quantities require global continuation analysis, which changes the source-duopoly best-response geometry and equilibrium correspondence.  Downstream exit is essential for the positive-slope response piece, but not for multiplicity itself: cap-induced asymmetric equilibria can already coexist with the symmetric equilibrium while both downstream firms remain active at every feasible sourcing history.  The exact thresholds and multiplicity results remain tied to the Shy--Stenbacka functional form; extending them to broader demand or monitoring-cost classes would require separate analysis.
